\Qhead{4.6}{Волноводы, длинные линии и резонаторы}{Волноводы, длинные линии и резонаторы. Критическая частота и критическая длина волновода. TE-, TH- и TEM-волны. Диэлектрические волноводы. Периодические структуры и замедляющие системы. Волновое сопротивление.}
\Subfile{4.6.1}{sections/04_generation/q06/s01}
\Subfile{4.6.2}{sections/04_generation/q06/s02}

\recap{Линия описывается распределёнными законами Кирхгофа; согласование устраняет отражение. Граничные условия создают резонатор. Периодические структуры замедляют фазу нужной гармоники до скорости электронов.}
\board{Элемент линии, телеграфные уравнения, $Z_0$, $\Gamma$; TE/TM/TEM; диаграмма синхронизма пучка.}
\oral{Почему фазовая скорость обычного пустого волновода неудобна для синхронного взаимодействия с медленными электронами?}
\bait{Можно оставить четвертьволновое преобразование: оно следует из противоположных узлов напряжения и тока стоячей волны.}
\sources{Биленко и соавт.~\cite{bilenko}, §2.3.1, PDF-с.~76--82; Белокопытов~\cite{belokopytov1}, §2 и §4. Соглашение фазоров для входного импеданса явно оговорено.}
